\section{Discussion}
\label{S:discussion}

\subsection{What the stage comparison establishes}

The principal finding is a measured allocation trade-off. On the fixed SAL prediction artifact,
structural localization removes many non-target sentences with little loss of labelled positives.
The detector then reduces candidate load, and selective verification improves precision while
sacrificing some of the detector's recall. The RCT-prompt $F_1$ increase of 0.054 is useful for
this configuration, but its 13.3\% loss of detector true positives is also material. Verification
is therefore an optional operating-point choice, not a universally beneficial final step.

The result depends on where a stage is evaluated. A full-article detector sees a low-prevalence
stream, whereas curated conclusion sentences or a balanced recognition sample present a different
false-positive burden. Changing the negative sample changes precision and $F_1$ even when positive
recall is unchanged. Published values must consequently be matched by label definition, candidate
pool, prevalence, and metric before being treated as head-to-head comparisons.

\subsection{How the findings answer the research questions}

For RQ1, the stage counts explain both the gain and its cost. Localization retains almost all
labelled positives, whereas most of the final recall loss occurs in detection and verification.
A practitioner prioritizing discovery can retain the detector's broader candidate set; a reader
prioritizing a shorter review list may prefer the verified output. Both outputs preserve the
original sentence and can be exposed with their decision provenance. The experiment supports
this operating-point distinction rather than a single mandatory setting for every application.

For RQ2, the recorded SAL comparison favors the fine-tuned biomedical detector among the eight
saved vectors, while the separate full-stream baseline grid also shows a strong lexical baseline.
The breadth results caution against assuming that its ranking or verifier policy transfers unchanged.
Future-work recognition, own-study limitation recognition, and broader research-gap recognition
require different judgments. A reusable system should therefore expose the target policy and
its corresponding evaluation. The conflicting suffix in the current gap-matched experiment
also shows that policy must be checked in the assembled request, not just in the base prompt.
The present evidence is strongest for the saved SAL configuration and weaker for aggregate-only
transfer studies and silver-label internal observations.

For RQ3, selective processing provides a clear reduction in estimated verifier input under the
specified accounting profile. This result explains why the staged design is attractive when
many articles must be inspected. It does not determine the best overall deployment architecture:
a small local detector adds its own preparation and maintenance costs, and a hosted model's
response reliability affects both effectiveness and retries. These quantities should be measured
when making a system-level efficiency claim. The branch's projected cross-paper batching, prompt
caching and long-paper savings are scenarios rather than measured deployment improvements;
we do not promote them to billed-cost or throughput results.

\subsection{Attribution policy and reference labels}

The Sci-Sentence development experiments illustrate the importance of matching the intended
prompt policy to the reference task, although the conflicting shared wrapper prevents a clean
causal test of that change. A verifier instructed to reject criticism of prior studies can correctly follow that
instruction while disagreeing with a dataset that labels those criticisms as Research Gap or
Limitation. This is a target mismatch, not evidence that either definition is intrinsically wrong.
Likewise, a statement about a study's methodological constraint and an adjacent consequence can
be scored differently under sentence- or span-oriented guidelines.

Inspection of disagreements can reveal annotation ambiguities, missed labels, model errors,
and sentence-boundary errors. It cannot by itself establish that most false positives are correct
or that an LLM achieves a much higher ``true'' precision than its measured precision. Such a claim
requires independent, preferably domain-expert, adjudication of a representative disagreement
sample. We therefore retain the original benchmark labels in all reported metrics. Neighbor-credit
scores, if used, should be interpreted as positional sensitivity rather than corrected gold.

Source grounding has similar limits. Returning an existing parsed sentence avoids inventing the
sentence's wording, but does not establish that it is a genuine author-stated gap, that its subject
is the authors' own work, or that the gap remains unresolved. Future downstream systems would need
separate attribution, resolution, and novelty checks before using these extractions to propose research.

\subsection{Efficiency and deployment}

Selecting candidate sentences substantially reduces estimated LLM input in the measured SAL
accounting profile. The 3.74-fold ratio incorporates some repeated-prompt and context overhead,
which makes it more informative than simply quoting the percentage of sentences removed.
However, a complete system cost also includes parsing, embeddings, detector training and inference,
response retries, and output tokens. API billing depends on the actual tokenizer and provider
pricing. Prompt caching and cross-paper batching could change that profile, but their projected
savings are not measured deployment results here.

The ten-article full-LLM comparison is encouraging but limited. It uses one economical model,
one selected sample, a different extraction prompt, and a baseline with response-format failures.
A stronger comparison would hold the annotation policy fixed, save all responses, use equivalent
retry/chunking budgets, and measure both stage-wise effectiveness and actual resource use.
The present result should guide that evaluation rather than substitute for it.

\subsection{Implications for evidence-based research support}

An extracted statement is useful as a traceable starting point for literature analysis. It can
help readers locate a paper's qualifications, collect proposed follow-up studies, and compare
how a topic is framed across articles. Keeping the source wording and identifier allows those
readers to inspect the surrounding argument rather than relying on an isolated model summary.
However, aggregation requires further decisions: two sentences can express the same issue,
a proposed study may already have been completed, and a frequently mentioned limitation may
reflect reporting convention rather than scientific priority.

Broader human evaluation and venue-corpus analysis can address some of these questions. Such
analyses must connect their article samples, annotation definitions, and run configurations to
the extraction protocol. Annotation agreement, venue-level statement frequencies, and the
usefulness of generated research ideas each require their own measurements. They cannot be
inferred from a sentence-detection score or a reduction in verifier input alone.

\subsection{Threats to validity}

\paragraph{Selection and sampling.}
The verifier prompts were refined after inspecting benchmark feedback. Descriptive bootstrap
intervals capture variation among the saved articles, not the optimism introduced by that
selection. The next evaluation should freeze prompts, head, threshold, context, and batching before
examining an independent article set. The cost subset is additionally conditioned on at least two
labelled limitations per article; its outcomes need not describe papers with few or no such statements.
The FWS breadth selection excludes exact training sentences but does not certify exclusion of
all related or same-paper training material.

\paragraph{Measurement and transfer.}
Recovered BioC tags simplify full-article localization relative to noisy PDFs. Curated FWS and
Sci-Sentence inputs do not test that front end at all. The internal arXiv references and the larger documented reference are model-labelled.
The 362-row interface has no completed human results and samples biomedical disagreements,
so it cannot establish reference accuracy across both domains. Inter-model kappa is not a
substitute for human agreement. The small arXiv reference is silver,
positive-only, and contains an expanded label ontology; it supports limited target-specific
recall observations, not population accuracy or human agreement. Proxy training-label mappings,
including Extension to future work, further limit interpretation of cross-domain rankings.

\paragraph{Reproducibility.}
Principal SAL decisions survive and their scores can be checked without an API, but the fine-tuned detectors
still lack checkpoints and full training provenance; the newly supplied frozen-head regenerator
resolves only part of that gap. Some experiment scripts retain machine-specific paths or imports
of helpers absent from the minimal branch. Experiment base prompts have been centralized, but the historic RCT alias still needs to be
archived with its exact resolved wording. Models referenced through a provider's \texttt{latest}
alias may change, and temperature zero does not establish repeatability across versions or batch
compositions. Aggregate-only breadth and full-article results cannot support additional paired
significance analyses or error adjudication without their original predictions and responses.

These limits bound the current conclusions: a particular saved full-article funnel improves
its recorded precision and $F_1$, and selective processing reduces estimated input volume.
Independent human-annotated results, broader venue findings, and evaluated research-idea quality
are outside the evidence reported in this study.
